\subsection{An arbitrary smooth root with full Hopf subtrees}

\begin{statement}{C.2.UN.61}{Proved}{The exact rooted Hopf-tree criterion and compact occurrence}
Let $S$ be a specified compact connected reduced SNC surface whose
incidence graph is a nontrivial tree. Distinguish one smooth component
$P$, with no initial restriction on its geometric genus or surface
type, and suppose that all other components are primary or secondary
Hopf surfaces. Retain the actual elliptic gluing maps. The following
finite geometric test is necessary and sufficient for $S$ to occur
as a full fibre of a smooth-total degeneration of connected complex
two-tori.

\begin{enumerate}
\item The actual normal lines at every double curve are dual.
\item The root $P$ has a fixed-Jacobian elliptic quasi-bundle
$q_P:P\to B$ over a smooth elliptic curve, with identity integral
monodromy and only smooth multiple supports. Its marked boundary
curves are fibre supports, and every multiple fibre is marked.
Every other component is either a nonlinear algebraic-dimension-zero
Hopf leaf or an elliptic Hopf quasi-bundle over $\PP^1$ whose
multiple fibres are all marked double curves.
\item For each edge $e$, let $m_e$ be the exact finite order of
its normal line. It is the full elliptic multiplicity at each
elliptic endpoint and the resonance at a nonlinear endpoint.
Set $h_P=0$ and
\[
                  h_v=\sum_{e\in[P,v]}m_e.
                  \tag{C.PT.1}
\]
There is an integer $R>0$ such that every nonlinear leaf has
$h_v=R$ and every elliptic Hopf vertex has $h_v<R$.
In the absence of nonlinear leaves any $R>\max_vh_v$ passes.
\item Perform the exact leaf-pruning operations of C.2.UN.60,
retaining the complete inverse logarithmic charts, primitive fibre
directions, secondary integers and torsor transitions. The final
surface is a principal elliptic torsor $T_0\to B$. Its integral
Chern vector in $H^2(B,\Lambda_F)$ is zero.
\end{enumerate}

The normal-root maps identify the one fixed Jacobian $F$ along
the tree. The last condition refers to the actual analytic
inverse torsor; it does not replace it by its rational Chern
data. The pruning result is independent, up to the actual
marked isomorphisms of UN60, of the order of disjoint operations.

Every support passing the test has a realization over $\PP^1$
with smooth compact total threefold. Its only singular fibre
is the given reduced $S$; every other fibre is a smooth
two-torus. All reduced fibre supports have Euler number zero.
Every component and every specified elliptic gluing map is
retained.

In every actual occurrence the root is minimal,
$h^q(P,\OO_P)=h^q(S,\OO_S)=(1,2,1)$, and its actual canonical
divisor is
\[
                 K_P=\sum_{e\ni P}(m_e-1)J_e.
                 \tag{C.PT.2}
\]
It is the unique minimum-order component for canonical evaluation;
the differences of the other orders are exactly (C.PT.1).
The zero-edge case, excluded from the nontrivial tree here,
is the smooth torus fibre already treated in C.2.UN.5.
\end{statement}

\begin{proof}
\par\medskip\noindent\textbf{The root cohomology is forced by the actual incidence maps.}
Every double curve lies in a Hopf component, so it is smooth
elliptic by the curve argument of UN54. Distinct such curves
on a Hopf component are disjoint, so there are no triple
points. In particular the curves $J_e$ on $P$ are also
pairwise disjoint. The exact normalization sequence is
\[
 0\longrightarrow\OO_S
 \longrightarrow\OO_P\oplus\bigoplus_{v\ne P}\OO_{H_v}
 \longrightarrow\bigoplus_e\OO_{J_e}\longrightarrow0.
 \tag{C.PT.3}
\]
Every Hopf restriction $H^1(\OO_{H_v})\to H^1(\OO_{J_e})$
is nonzero, with its actual gluing map, as proved in UN54.
Root the tree at $P$ and order the Hopf columns from the
leaves inward. The submatrix of the $H^1$ incidence map
on those columns is square, with one column and row per
nonroot vertex. Elimination at a leaf has its nonzero
Hopf restriction as pivot. Repeating gives an invertible
matrix. The full $H^1$ incidence map is therefore onto,
regardless of the root restriction maps.

The $H^0$ incidence map of a tree is also onto. Since
$H^2(\OO_H)=0$ on a Hopf surface, the actual long exact
sequence of (C.PT.3) gives
\[
 H^2(S,\OO_S)\simeq H^2(P,\OO_P),\qquad
                 h^1(S,\OO_S)=h^1(P,\OO_P).
 \tag{C.PT.4}
\]
The smooth total map is flat over the disc. Upper
semicontinuity gives $h^2(S,\OO_S)\ge1$, while the
compact-core bound UN28 gives $p_g(P)\le1$. Thus
$p_g(P)=1$. Coherent Euler characteristic in this proper
flat family is zero. Additivity in (C.PT.3), since Hopf
surfaces and elliptic curves have coherent Euler zero,
then gives $\chi(\OO_P)=0$ and $q(P)=2$. This proves
the asserted coherent dimensions without an initial
geometric-genus hypothesis on $P$.

\par\medskip\noindent\textbf{The actual dualizing form and the unique minimum.}
Use the residue sequence (C.HR.7) at every Hopf vertex.
At a terminal Hopf vertex its one nonzero restriction
coefficient forces the sole residue of a global
dualizing form to be zero. Since a Hopf surface has no
holomorphic two-form, the form on that component is zero.
Remove that leaf and repeat. All Hopf forms and every
edge residue are zero. The remaining form on $P$ is
an ordinary holomorphic two-form. Conversely such a form,
with zero on all Hopf components, glues by the zero-residue
rule. These are inverse maps, giving the actual isomorphism
\[
                  H^0(S,\omega_S)\simeq H^0(P,K_P).
 \tag{C.PT.5}
\]

Take a canonical evaluation generator, with vertical orders
$k_v$, and divide it by $t^{\min k_v}$. Its nonzero restriction
to the reduced fibre spans (C.PT.5). The local adjunction map
(C.HR.8) identifies its nonzero components precisely with the
minimum-order components. Hence the minimum is attained
only at $P$. Writing $k_P=k_0$, the same actual residue
calculation on $P$ gives a holomorphic two-form with divisor
\[
          K_P=\sum_{e=(P,w)}(k_w-k_0-1)J_e,\qquad
                          k_w>k_0.
 \tag{C.PT.6}
\]
All coefficients are nonnegative integers; there is no
additional zero away from the boundary, because the
evaluation has only its vertical component divisors.
If $E$ were a $(-1)$-curve, adjunction would give
$K_P\cdot E=-1$. The effective divisor in (C.PT.6)
is supported on elliptic curves, none equal to $E$,
and has nonnegative intersection with $E$. This
contradiction proves that $P$ is minimal.

\par\medskip\noindent\textbf{The actual elliptic fibration on the root.}
Its nonzero canonical form gives $\kappa(P)\ge0$.
Every component has full multiplicity one, and the
given total space and central support are already
smooth and SNC. Thus the actual common-root map in
UN31 has degree one and its component over $P$ is
$P$ itself. That theorem makes $P$ a two-torus or
gives on it a fixed-modulus, identity-monodromy
elliptic quasi-bundle with only smooth multiples.
Its base has genus $p_g(P)=1$.

We record an exact principal-torsor comparison used twice
below. For a principal $F=\C/\Lambda_F$-torsor over a
smooth elliptic curve $B$, the sequence
\[
 0\longrightarrow\Lambda_F
 \longrightarrow\OO_B\longrightarrow F_B
 \longrightarrow0
 \tag{C.PT.7}
\]
maps its actual analytic class to its integral Chern
vector. If that vector is zero, the class comes from
$H^1(B,\OO_B)$. On $B=\C/\Lambda_B$, a representative
$a\,d\overline z$ has additive transition functions
$a\overline\gamma$ for $\gamma\in\Lambda_B$.
Its actual total torsor is therefore
\[
 \C^2/
 \left\langle
  (z,w)\mapsto(z+\gamma,w+a\overline\gamma),\
  (z,w)\mapsto(z,w+\lambda):
   \gamma\in\Lambda_B,\ \lambda\in\Lambda_F
 \right\rangle .
 \tag{C.PT.8}
\]
The four real periods are independent by projection
to $\Lambda_B$ and then to $\Lambda_F$; they give
a discrete cocompact lattice. This is a complex
two-torus, with the original torsor class represented
by the same additive transitions. A different lift
through (C.PT.7) gives its isomorphic original
$F$-torsor. No section of the original torsor has
been assumed.

We now identify every marked elliptic curve on $P$
with a fibre of one such elliptic fibration.
On a torus an elliptic curve is a translate of an
elliptic subtorus: lifting its map to universal
covers makes its derivative a periodic holomorphic
function, hence constant. Two different elliptic
directions give an isogeny from their product to
the two-torus, so their translates intersect.
The pairwise disjoint boundary curves therefore
have the same direction and are fibres of the
actual quotient by that elliptic subtorus.

In the other alternative let $q:P\to B$ be the
fibration supplied by UN31. If it has a multiple
fibre, any nonvertical marked elliptic curve would
map finitely to the elliptic $B$. Riemann--Hurwitz
makes that map unramified. But at a multiple fibre
the actual local function is $s^m$, $m>1$,
so the restriction to that curve has positive
ramification, a contradiction.
If there is no multiple fibre, $q$ is a principal
$F$-torsor. A nonvertical elliptic curve would give
a section after its finite base pullback. Pullback
on $H^2(B,\Lambda_F)$ multiplies its Chern vector
by the positive covering degree. The vector group
is torsion-free, so the original vector would be
zero. Formula (C.PT.8) then makes $P$ a torus,
already handled above. Otherwise all marked
curves are vertical. Since every reduced fibre
is one smooth elliptic curve, each is exactly
a marked fibre support. This proves the required
actual fibration $q_P$ in all cases.

\par\medskip\noindent\textbf{Finite normal roots and the exact height equations.}
The normal-root chart (C.HT.4) gives each root
boundary normal a finite exact order equal to its
full elliptic multiplicity. Normal duality carries
that order to the adjacent Hopf component.
An algebraic-dimension-zero diagonal Hopf has only
nontorsion axis normals, so it cannot be adjacent
to $P$. At an elliptic Hopf all curve normals have
finite order, and a nonlinear Hopf has only its
single torsion-normal axis. Propagating along the
tree excludes every algebraic-dimension-zero
diagonal vertex. Thus every other vertex is an
elliptic Hopf or a nonlinear leaf.

For each elliptic component, the actual root
charts give the global meromorphic form
$q^*\theta\wedge dz$ and its exact order
$m\operatorname{ord}\theta+m-1$, as in
(C.HT.5)--(C.HT.6). The canonical residue form
on that component has divisor only on its
marked curves. At an unmarked multiple fibre
its order would be zero, which is incompatible
with $m\operatorname{ord}\theta+m-1=0$ for
$m>1$. This proves the absence of all unmarked
multiples on $P$ and on every elliptic Hopf
vertex.

The dual-normal maps identify the connected
normal-root elliptic covers. Beginning with
the fixed Jacobian of $q_P$ and following the
unique path in the tree gives its actual
identification with every later Jacobian and
with each nonlinear axis root cover. In
particular the exact normal order $m_e$ is
the same at both endpoints.

At every elliptic Hopf vertex the residue
calculation gives the integer equation
\[
 \sum_{e=(v,w)}
       \frac{k_w-k_v}{m_e}
                        =\operatorname{val}(v)-2.
 \tag{C.PT.9}
\]
A terminal elliptic vertex gives the outward
increment $m_e$. A terminal nonlinear vertex
gives the same increment from its exact
axis form $dz\wedge dw/w^{m_e+1}$ and its
infinite-order effective axis line.
The finite leaf induction of UN58 applies
without an equation at the root: once every
child term is one, (C.PT.9) makes its parent
term minus one. Therefore every outward
increment is $m_e$, proving (C.PT.1).
Substituting the root increments into
(C.PT.6) proves (C.PT.2).

The actual root collars $xy=t$ and the full
secondary boundary matrices of UN50 identify
every regular elliptic fibre and every
nonlinear real core with one transported
integral two-cycle $\gamma_t$ in the nearby
torus. They are exactly the maps
(C.HT.10)--(C.HT.12); the genus of the root
base does not change those collar maps.
For the canonical evaluation divided by
$t^{k_0}$, its leading form on an elliptic
vertex restricts to zero on the elliptic
curve. Its period is therefore
$O(t^{h_v+1})$, including at $P$.
At a nonlinear leaf it is
\[
 t^{h_v}
       \left(\frac{\pm2\pi i c\beta^{-m_e}}{n}
                         +O(t)\right),\qquad c\ne0.
 \tag{C.PT.10}
\]
These are the same holomorphic period, by the
actual integral collar identifications.
One nonlinear leaf fixes its finite positive
order $R$; all nonlinear heights equal $R$,
and every elliptic vertex has smaller height.
This proves the stated height test.

\par\medskip\noindent\textbf{The exact inverse torsor and its integral obstruction.}
Perform the operations (C.PR.6) from the leaves
inward. After its children have been removed,
an elliptic Hopf vertex is the actual one-mark
torsor of (C.PR.1)--(C.PR.5), hence is a
recorded diagonal Hopf pair or the product
ruled cap. Nonlinear leaves use the same
inverse product map, with their full
secondary integer and primitive fibre
direction. The root always remains its
fixed-Jacobian quasi-bundle over the original
elliptic $B$. Each operation replaces exactly
one marked multiple-fibre tube by its
displayed ordinary product tube; a product
cap leaves that ordinary tube unchanged.
There are no unmarked multiples. After all
Hopf vertices have been removed, the final
root is therefore a principal $F$-torsor
$T_0\to B$, with its full analytic transition
class retained.

The exact boundary homomorphism of every
removed compact piece is (C.PR.9): it kills
the one meridian $\delta$ and retains the
two generators $e,h$. The inverse product
filling has those identical maps. Applying
the actual Van Kampen pushouts successively,
as in UN60, gives
\[
                         \pi_1(T_0)\simeq\pi_1(F_t)
                                             \simeq\Z^4.
 \tag{C.PT.11}
\]
Only these smooth collar maps are used for
an arbitrary original degeneration; no
holomorphic contraction of its neighbourhood
is asserted.

For a principal elliptic bundle over an
elliptic curve the homotopy sequence has
an injected fibre lattice, since
$\pi_2(B)=0$. The group is the central
extension
\[
 1\longrightarrow\Lambda_F
 \longrightarrow\pi_1(T_0)
 \longrightarrow\Lambda_B\longrightarrow1.
 \tag{C.PT.12}
\]
The commutator of lifts of an oriented
base basis is its integral Chern vector.
This follows directly by trivializing over
the usual base parallelogram: its boundary
commutator is the clutching loop of the
fibre torus. Equation (C.PT.11) makes that
commutator zero in the injected lattice.
Thus the Chern vector is zero. Formula
(C.PT.8) identifies this very inverse
torsor, including its analytic class,
with a complex two-torus. This proves
necessity of the last test. Independence
of disjoint pruning orders is the actual
chart comparison of UN60.

\par\medskip\noindent\textbf{Compact reconstruction of the complete prescribed tree.}
Assume all the tests pass and retain the
complete pruning record. Formula (C.PT.8)
gives the actual initial torus
$T_0=\C^2/\Gamma$. Its elliptic fibre
lattice is a primitive rank-two subgroup
of $\Gamma$, because the quotient base
lattice is free. Choose $R$ as in the
height test and put
\[
 V_t(w)=w+t^R J\overline w,\qquad
              J=\begin{pmatrix}0&1\\-1&0\end{pmatrix}.
 \tag{C.PT.13}
\]
For any real basis of $\Gamma$ its image
has determinant $(1+|t|^{2R})^2$ times
the original nonzero determinant.
These holomorphic period sections
therefore give a proper smooth torus
family on the finite chart. At infinity
the frame divided by $t^R$ gives
$J\overline w+s^Rw$, $s=t^{-1}$.
It is again real-linearly invertible
at every point, including $s=0$.
The two charts consequently give an
actual compact smooth torus family
over $\PP^1$, with central member
exactly $T_0$.

For an oriented primitive fibre basis
$\lambda,\mu=c\lambda$,
$\operatorname{Im}c>0$, its exact period
in this frame is
\[
 \det_\C(V_t\lambda,V_t\mu)
       =2i\operatorname{Im}(c)\|\lambda\|^2t^R.
 \tag{C.PT.14}
\]
The coefficient is nonzero. Near each
recorded central elliptic fibre use
the actual rank-two period cover and
the primitive-basis change (C.PR.13).
This gives
$(p,w)\mapsto(p+\delta(t),\gamma(t)w)$
with $\delta(t)=t^R d(t)$, $d(0)\ne0$,
and the original marked Jacobian and
primitive fibre direction.

Reverse the pruning record. At a
recorded Hopf step of weight $m$ use
the free weighted chart (C.PR.14),
including its full secondary integer.
Its child's affine pencil has the
exact displacement $\delta(t)/t^m$.
At a product-cap step use the ordinary
blowup (C.PR.16), with displacement
$\delta(t)/t$. Thus the remaining
order at every vertex is $R-h_v$.
The height test makes each recorded
elliptic child linear and every
recorded nonlinear terminal child
nonlinear, with its prescribed
resonance and contraction. The
recorded primitive basis changes
retain all linear gluing maps,
and the retained transition charts
restore all elliptic translations
and every analytic torsor.

All local compactness comparisons
are precisely (C.PR.17)--(C.PR.18).
Exponential contraction dominates
the polynomial exceptional norm
distortion, and the old base
displacement is
$O(|\delta(t)|\,|\log|x||)$.
Its cluster set lies in the old
marked elliptic curve. These exact
bounds give disjoint Hausdorff
proper replacements at each stage;
ordinary caps are proper blowups.
There are finitely many stages.
Each is the identity on the
punctured family and outside its
chosen small central neighbourhood.
They therefore preserve the
compact family over the rest of
$\PP^1$, including infinity.

Induction through the recorded
inverse charts restores precisely
the specified $P$, every Hopf
component and each prescribed
elliptic gluing map. Their normal
equations are the free coarse
$t=xW$ charts or the ordinary
$t=xy$ charts, so the full fibre
is reduced and the total space
is smooth. Every other fibre is
the unchanged smooth torus.
The root quasi-bundle, the Hopf
components and all double
elliptic curves have Euler
number zero. Additivity proves
the stated zero Euler number
for $S$ and for all other supports.

The rooted incidence matrices,
the integral Chern commutator,
the exact torsor period matrix
and the height equations are
checked in
\path{code/A_rooted_hopf_trees.py};
the primitive period and nested
order identities are also
certified in
\path{code/A_tree_periods.py}.
\end{proof}
